7.1 一体系由三个全同玻色子组成, 不考虑粒子间的相互作用. 已知可能的单粒子态为 $\varphi_{1}$ 与 $\varphi_{2}$ , 相应的能量为 $\varepsilon_{1}$ 与 $\varepsilon_{2}$ , 写出体系所有可能态的波函数和能量.

\vfill\newpage

7.2 两个自旋 $s = 1 / 2$ 的全同粒子在同一谐振子势场中运动， $V(r) = \frac{1}{2}\mu \omega^2 r^2$ ，不考虑两粒子之间的相互作用，求一粒子处于单粒子基态，另一粒子处于在 $z$ 方向运动的单粒子第一激发态的体系波函数和能量，并求体系总角动量量子数 $(j, m_j)$ 的可能值。

\vfill\newpage

7.3 (1) 在一维无限深方势阱 $V(x) = \begin{cases} 0, & 0 < x < a \\ \infty, & x < 0, x > a \end{cases}$ 中有两个自旋 $s = 0$ 的全同粒子, 粒子之间不存在相互作用, 写出体系最低两个能级, 指出简并度, 并给出相应的波函数; 

(2) 同(1), 粒子具有自旋 $s = 1/2$ ; 

(3) 同(2), 但粒子之间存在同自旋有关的相互作用力势 $V = A\hat{S}_1 \cdot \hat{S}_2 (A > 0)$ .

\vfill\newpage

7.4 设绝对零度时, 在三维各向同性谐振子势 $V(r) = \frac{1}{2}\mu \omega^2 r^2$ 中有 20 个自旋 $s = 1 / 2$ 质量为 $\mu$ 的全同粒子组成的体系. 忽略粒子之间的相互作用. 已知这 20 个粒子的平均能量为 $3\mathrm{eV}$ . 

(1) 如果同样的温度下该势场中有 12 个这样的粒子组成的体系, 其平均能量是什么? 

(2) 如果同样的温度下该势场中有 17 个自旋 $s = 0$ 质量仍为 $\mu$ 的全同粒子组成的体系, 其平均能量是什么?

\vfill\newpage

7.5 设有两个质量为 $\mu$ 自旋 $s = 1 / 2$ 的全同粒子, 在同一势场 $V = \frac{1}{2} kx^{2}$ 中作一维谐振子运动, 即 $V_{i} = \frac{1}{2} kx_{i}^{2}(i = 1,2)$ . 两个粒子之间的相互作用力势为 $\frac{1}{2}\alpha k(x_{1} - x_{2})^{2}, \alpha > -1 / 2$ . 试求体系的能级, 并指出哪些能级属于自旋单态? 哪些能级属于自旋三重态?

\vfill\newpage

7.6 假设氦原子中两个电子被没有自旋的玻色子取代(质量和电荷不变), 试讨论能级的变化.

\vfill\newpage

7.7 两个质量为 $\mu$ 的粒子处于边长为 $a > b > c$ 的立方体盒中, 粒子间相互作用势 $V = A\delta(r_1 - r_2)$ 可视为微扰. 在下列条件下, 用一级微扰方法计算体系的最低能量: (1) 粒子非全同; (2) 零自旋的全同粒子; (3) 自旋为 $1/2$ 的全同粒子, 并处于总自旋 $s = 1$ 的态上.

\vfill\newpage

7.8 假设两个质量为 $m_{\mathrm{q}} = 70 \, \mathrm{MeV} / c^2$ 的夸克可以通过相互作用位势 $V(r) = -a(\hat{\sigma}_1 \cdot \hat{\sigma}_2 - b)r^2$ 束缚在一起, 其中 $r$ 是两个夸克之间的距离, $b$ 是一个待定的参数, $a = 68.99 \, \mathrm{MeV} / \mathrm{fm}^2$ . 

(1) $b$ 取什么值才能使两个夸克束缚在一起? 

(2) 设两个夸克是不同类型的, 并取 $b = 3/2$ . 试求基态能量和简并度; 

(3) 设两个夸克是同一类型的, 并取 $b = 3/2$ . 试求基态能量和简并度; 

(4) 令 $b = 0$ , 求两个全同夸克在基态的方均根距离. 已知 $\hbar c = 197.3 \, \mathrm{MeV} \cdot \mathrm{fm}$ .

\vfill\newpage

7.9 一体系由两个质量为 $m$ 的一维全同粒子组成, 两粒子之间的相互作用为 $a(x_{1} - x_{2})^{2} / 2 (a > 0)$ . 

(1) 若粒子自旋为 0, 写出它们的相对运动态的能量和波函数; 

(2) 若粒子自旋 $s = 1 / 2$ , 写出它们的相对运动基态及第一激发态的能量和波函数.

\vfill\newpage

7.10 氯化钠晶体中有些负离子空穴. 每个空穴束缚一个电子, 可以认为, 这些电子被束缚在一个尺度为晶格常数的三维无限深势阱中, 晶体处于室温. 试粗略估计被这些电子强烈吸收的电磁波的最长波长. 已知 $\hbar c = 197\mathrm{MeV}\cdot \mathrm{fm}$ , 晶格常数 $a = 0.1\mathrm{nm} = 10^{5}\mathrm{fm}$ , 电子质量 $\mu c^2 = 0.511\mathrm{MeV}$ .

\vfill\newpage

7.11 设 $a^+$ 与 $a$ 是玻色子在某单粒子态上的产生算符与湮没算符, 满足对易关系 $[a, a^+] = 1$ . $\hat{N} = a^+ a$ 是该单粒子态上的粒子占有数算符, $\left|n\right\rangle$ 是 $\hat{N}$ 的本征值为 $n$ 的本征态. 

(1) 计算对易关系 $[a, \hat{N}], [a^+, \hat{N}]$ ; 

(2) 证明
\begin{equation}\tag{1}
    a | n \rangle = \sqrt {n} | n - 1 \rangle    
\end{equation}
\begin{equation}\tag{2}
    a ^ {+} | n \rangle = \sqrt {n + 1} | n + 1 \rangle
\end{equation}

(3) 求 $\hat{N}$ 的本征值 $n$ .

\vfill\newpage

7.12 算符 $a_{i}$ 与 $a_{j}^{+}(i,j = 1,2)$ 满足对易关系 $[a_i,a_j^+] = \delta_{ij}$ ， $[a_i,a_j] = 0$ ， $[a_i^+, a_j^+] = 0$ 。体系的哈密顿量为 $\hat{H} = \hbar \omega_0(a_1^+ a_1 + a_2^+ a_2) + i\hbar \omega_1(a_1^+ a_2 - a_2^+ a_1)$ ，其中 $\omega_0$ 与 $\omega_1$ 均为正实数，且 $\omega_1 \ll \omega_0$ 。试用微扰方法计算体系的第一激发态的一级近似能量，并同精确能量比较。

\vfill\newpage

7.13 $a$ 与 $a^{+}$ 是费米子体系的某个单粒子态的湮没与产生算符, 满足反对易关系 $\{a, a^{+}\} \equiv aa^{+} + a^{+}a = 1$ , $a^{2} = (a^{+})^{2} = 0$ . 以 $\hat{N} = a^{+}a$ 表示该单粒子态的粒子占有数算符, $\left|n\right\rangle$ 为 $\hat{N}$ 的本征值为 $n$ 的本征态矢. 

(1) 求 $\hat{N}$ 的本征值 $n$ ; 

(2) 计算对易关系式 $[\hat{N}, a], [\hat{N}, a^{+}]$ ; 

(3) 证明 $a\left|n\right\rangle = \sqrt{n}\left|n - 1\right\rangle$ , $a^{+}\left|n\right\rangle = \sqrt{1 - n}\left|n + 1\right\rangle$ ; 

(4) 在 $\hat{N}$ 表象求 $a$ 与 $a^{+}$ 的矩阵表示.

\vfill\newpage

7.14 算符 $a_{i}$ 与 $a_{j}^{+}$ $(i,j = 1,2)$ 满足反对易关系 $\{a_i,a_j\} = \{a_i^+,a_j^+\} = 0$ ， $\{a_i,a_j^+\} \equiv a_i a_j^+ + a_j^+ a_i = \delta_{ij}$ ，

(1)试求哈密顿量 $\hat{H}_0 = \hbar \omega_1 a_1^+ a_1 + \hbar \omega_2 a_2^+ a_2$ ( $\omega_2 > \omega_1 > 0$ ) 的能谱和本征态矢；

(2)在 $\hat{H}_0$ 表象给出算符 $\hat{Q} = a_1 a_2$ 和 $\hat{W} = a_1^+ a_2^+$ 的矩阵表示；

(3)设 $\hat{H}' = \varepsilon (a_1^+ a_2^+ - a_1 a_2)$ 为微扰，求 $\hat{H}_0$ 的基态在计入微扰后的二级近似能量和一级近似态矢.

\vfill\newpage



7.15 设 $a$ 和 $a^{+}$ 是湮没算符和产生算符, 满足对易关系 $[a, a^{+}] = 1$ . 体系的哈密顿量为
\begin{equation}\tag{1}
    \hat {H} = A a ^ {2} + B (a ^ {+}) ^ {2} + C a ^ {+} a + D
\end{equation}
请问 A, B, C, D 要满足什么条件, $\hat{H}$ 才是厄米算符? 求出体系的能量.

\vfill\newpage

7.16 某体系哈密顿量 $\hat{H} = \frac{5}{3} a^{+}a + \frac{2}{3} [a^{2} + (a^{+})^{2}]$ ，其中 $a = \frac{1}{\sqrt{2}} (Q + \mathrm{i}\hat{P})$ ， $a^{+} = \frac{1}{\sqrt{2}} (Q - \mathrm{i}\hat{P})$ ， $\hat{P}$ 与 $Q$ 满足基本对易关系 $[Q, \hat{P}] = \mathrm{i}$ ， $\hat{P} = -\mathrm{i}\frac{\mathrm{d}}{\mathrm{d}Q}$ . 试求 $\hat{H}$ 的本征值与基态波函数 $\psi_{0}(Q)$ .

\vfill\newpage

7.17 设哈密顿算符 $\hat{H} = \lambda a^{+}a + \varepsilon (a^{+} + a)$ , 其中 $\lambda$ 是正实数, $\varepsilon$ 是小参数, $a^{+}$ 与 $a$ 是玻色子产生算符与湮没算符. 求 $\hat{H}$ 的基态能量本征值(准至 $\varepsilon^2$ 级), 并同严格值比较.

\vfill\newpage

7.18 在粒子数表象中, 一维谐振子基态态矢 $|0\rangle$ 满足性质 $a|0\rangle = 0$ , 其中 $a$ 为湮没算符, $a = \sqrt{\frac{\mu\omega}{2\hbar}}\left(x + \frac{\mathrm{i}}{\mu\omega}\hat{p}\right)$ . 试利用此性质求出基态在动量表象中的波函数显示式 $\langle p|0\rangle = \psi_0(p)$ .

\vfill\newpage

7.19 一维谐振子哈密顿量 $\hat{H} = \hbar \omega\left(a^{+}a + \frac{1}{2}\right)$ ，且 $[a, a^{+}] = 1$ 。 

(1) 若 $|0\rangle$ 是归一化的基态态矢 ( $a|0\rangle = 0$ )，则第 $n$ 个激发态态矢为 $|n\rangle = N_{n}(a^{+})^{n}|0\rangle$ 。求归一化因子 $N_{n}$ 。

(2) 若外加一微扰 $\hat{H}' = ga^{+}a^{+}aa$ ，求第 $n$ 个激发态的能量本征值 (准至 $g$ 一级).

\vfill\newpage

7.20 一维谐振子受到微扰 $\hat{H}' = cx^2$ 的作用, 其中 $c$ 是常数. 在粒子数表象中 $x = \sqrt{\frac{\hbar}{2\mu\omega}} (a + a^{+})$ , $a$ 与 $a^{+}$ 分别是湮没算符和产生算符, 满足如下公式: $a|n\rangle = \sqrt{n}|n - 1\rangle$ , $a^{+}|n\rangle = \sqrt{n + 1}|n + 1\rangle$ , 其中 $|n\rangle$ 是一维谐振子哈密顿量 $\hat{H} = \hbar\omega\left(a^{+}a + \frac{1}{2}\right)$ 的本征态. 

(1)用微扰论, 准确到二级近似, 求能量修正值; 

(2)求能量的准确值, 并与微扰论给出的结果比较.

\vfill\newpage

7.21 考虑两个自旋 $s = 1 / 2$ ，质量为 $\mu$ 的全同粒子在三维空间内运动。假定两粒子的总动量为 0，两粒子间的相互作用势为 $V(r_1, r_2) = \frac{g}{|r_1 - r_2|} \sigma_1 \cdot \sigma_2$ ，其中 $g$ 是一个正实数， $\sigma_1$ 与 $\sigma_2$ 为两个粒子的泡利矩阵。求出两个粒子能形成的所有束缚态能级和相应的简并度。

\vfill\newpage

7.22 两个无相互作用的粒子置于一维无限深方势阱 $(0 < x < a)$ 中. 对于以下两种情况，写出两粒子体系可具有的两个最低能量值，相应的简并度，以及上述能级对应的所有二粒子波函数：

(1)两个自旋为 $1/2$ 的可区分粒子；

(2)两个自旋为 $1/2$ 的全同粒子.

\vfill\newpage

7.23 两个无相对作用粒子具有相同质量 $m$ ，在宽为 $a$ 的一维无限深方势阱中运动。

(1)写出体系4个最低能级的能量。

(2)对下述情况，分别求出体系4个最低能级的简并度：

(a)自旋为 $1/2$ 的全同粒子；

(b)自旋为 $1/2$ 的非全同粒子；

(c)自旋为1的全同粒子。

\vfill\newpage

7.24 求两个自旋 $s = 1 / 2$ 的关在一维无限深势阱 $V(x) = \left\{ \begin{array}{ll}0, & 0 < x < a\\ \infty , & x < 0,x > a \end{array} \right.$ 中，并以接触势 $U(x_{1},x_{2}) = c\delta (x_{1} - x_{2})(c\ll 1)$ 为相互作用的全同粒子系统的零级近似归一化波函数(考虑自旋态)，以接触势为微扰，求准确到 $c$ 的一次方的基态能量.

\vfill\newpage

7.25 考虑两个具有相同角频率 $\omega_0$ 的振子，哈密顿量为 $\hat{H}_1 = \hbar \omega_0 a_1^+ a_1$ ， $\hat{H}_2 = \hbar \omega_0 a_2^+ a_2$ 。记 $\hat{H}_1, \hat{H}_2$ 相应于本征值 $n_1 \hbar \omega_0$ 和 $n_2 \hbar \omega_0$ 的本征态为 $\left| n_1 n_2 \right\rangle$ ，零点能已略去。在两个振子具有相互作用后，其哈密顿量为
$$
\begin{array}{r l} & {\hat {H} = \hbar \omega_ {0} a _ {1} ^ {+} a _ {1} + \hbar \omega_ {0} a _ {2} ^ {+} a _ {2} + g a _ {1} ^ {+} a _ {2} + g a _ {2} ^ {+} a _ {1}} \\ & {\qquad = \left(a _ {1} ^ {+}, a _ {2} ^ {+}\right) \left( \begin{array}{c c} {\hbar \omega_ {0}} & {g} \\ {g} & {\hbar \omega_ {0}} \end{array} \right) \binom {a _ {1}} {a _ {2}}} \end{array}
$$
其中 g 为正实数. 因为有相互作用, $\left|n_{1}n_{2}\right\rangle$ 不是 $\hat{H}$ 的本征态. 求 $\hat{H}$ 的本征值. [提示: 使矩阵 $\begin{pmatrix}\hbar\omega_{0}&g\\g&\hbar\omega_{0}\end{pmatrix}$ 对角化]

\vfill\newpage

7.26 质量为 $\mu$ 的粒子处于三维各向同性谐振子势 $V(r) = \frac{1}{2}\mu \omega^2 r^2$ 中.

(1)求粒子的本征能量和相应的简并度；

(2)如再加上微扰 $\hat{H}' = bxz$ 的作用（ $b$ 为小的正实数），求粒子的基态和第一激发态能量的一级修正；

(3)如在以上势阱(含微扰)中，放入5个全同无相互作用的自旋 $s = 0$ 的粒子，求体系基态能量；

(4)如在以上势阱(含微扰)中，放入5个全同无相互作用的自旋 $s = 1 / 2$ 的粒子，求体系基态能量.

\vfill\newpage

7.27 $a_{i}^{+}, a_{i}$ 是玻色子在单粒子态 $\phi_{i}(i = 1,2)$ 上的产生算符与湮没算符，满足对易关系 $[a_{i}, a_{j}] = [a_{i}^{+}, a_{j}^{+}] = 0, [a_{i}, a_{j}^{+}] = \delta_{ij}$ . 令

$$
\hat {J} _ {x} = \frac {\hbar}{2} \left(a _ {1} ^ {+} a _ {2} + a _ {2} ^ {+} a _ {1}\right), \hat {J} _ {y} = \frac {\mathrm{i} \hbar}{2} \left(a _ {1} ^ {+} a _ {2} - a _ {2} ^ {+} a _ {1}\right), \hat {J} _ {z} = \frac {\hbar}{2} \left(a _ {2} ^ {+} a _ {2} - a _ {1} ^ {+} a _ {1}\right)
$$

(1) 证明

$$
\left[ \hat {J} _ {x}, \hat {J} _ {y} \right] = \mathrm{i} \hbar \hat {J} _ {z}, \left[ \hat {J} _ {y}, \hat {J} _ {z} \right] = \mathrm{i} \hbar \hat {J} _ {x}, \left[ \hat {J} _ {z}, \hat {J} _ {x} \right] = \mathrm{i} \hbar \hat {J} _ {y}
$$

(2) 证明

$$
\hat {J} ^ {2} = \hat {J} _ {x} ^ {2} + \hat {J} _ {y} ^ {2} + \hat {J} _ {z} ^ {2} = \hbar^ {2} \frac {\hat {N} _ {1} + \hat {N} _ {2}}{2} \left(\frac {\hat {N} _ {1} + \hat {N} _ {2}}{2} + 1\right)
$$
$$
\left[ \hat {J} ^ {2}, \hat {J} _ {\alpha} \right] = 0, \alpha = x, y, z
$$
其中 $\hat{N}_1 = a_1^+ a_1$ 与 $\hat{N}_2 = a_2^+ a_2$ 分别是 $\phi_1$ 与 $\phi_2$ 态上粒子占有数算符；

(3) 求 $\hat{J}^2$ 与 $\hat{J}_z$ 的共同本征态，以及它们的本征值.

\vfill\newpage